%\section{Public API Reference}

%The principal entry point is \code{CUDAHeat(source, variable=...)}. The \code{run()} method computes requested fields, whereas \code{detect()} returns an \code{EventDetectionResult}. Convenience methods expose the smoothed climatology, aligned threshold and in-memory exceedance mask. The result object supplies event selection, point-level catalogues, point time series, event-centred time series and lazy dense reconstruction. Numerical diagnostics are available through \code{validate()}, which compares CPU reference fields and GPU fields using the same temporal formulation.
